\documentclass[11pt,a4paper]{article}

% ------------------------------------------------
% Packages
% ------------------------------------------------
\usepackage[utf8]{inputenc}
\usepackage[T1]{fontenc}
\usepackage{lmodern}
\usepackage[a4paper,margin=2.2cm]{geometry}
\usepackage{microtype}
\usepackage{xcolor}
\usepackage{hyperref}
\usepackage{booktabs}
\usepackage{amsmath}
\usepackage{listings}
\usepackage{enumitem}

% ------------------------------------------------
% Hyperlinks and colors
% ------------------------------------------------
\definecolor{ectexblue}{RGB}{35,85,145}
\definecolor{codegray}{RGB}{245,245,245}

\hypersetup{
    colorlinks=true,
    linkcolor=ectexblue,
    citecolor=ectexblue,
    urlcolor=ectexblue,
    pdftitle={EcTeX: LaTeX for Eclipse},
    pdfauthor={EcTeX Project}
}

% ------------------------------------------------
% Code listings
% ------------------------------------------------
\lstset{
    basicstyle=\ttfamily\small,
    backgroundcolor=\color{codegray},
    frame=single,
    breaklines=true,
    columns=fullflexible,
    showstringspaces=false
}

% ------------------------------------------------
% Document information
% ------------------------------------------------
\title{\textbf{EcTeX: LaTeX for Eclipse}\\
\large An Integrated Editing and PDF Preview Environment}

\author{EcTeX Project}
\date{Version 1.1.0 -- 2026}

\begin{document}

\maketitle

\begin{abstract}
EcTeX is an Eclipse-based environment for writing and compiling
LaTeX documents. It integrates a syntax-highlighted text editor,
PDF preview, build commands and source-to-PDF synchronization
into a single workspace. By combining the capabilities of an
existing TeX distribution with the extensibility of Eclipse,
EcTeX aims to make scientific and technical document preparation
more convenient. This document describes its main features,
technical requirements and typical workflow.
\end{abstract}

\tableofcontents

\section{Introduction}

LaTeX is widely used for preparing scientific papers,
technical reports and other documents that require
structured text, mathematical notation and bibliographies.
Its separation of content and presentation makes it
particularly suitable for documents that need consistent
formatting and reproducible output.

However, working with LaTeX often involves switching
between a text editor, a terminal, a PDF viewer and
bibliography tools. EcTeX addresses this workflow by
integrating these activities into the Eclipse IDE.

The project provides a dedicated editor for \texttt{.tex}
files, an embedded PDF preview and convenient commands
for compiling documents and navigating between source
code and the generated PDF.

\section{Features}

\subsection{LaTeX editor}

EcTeX provides a specialized editor for \texttt{.tex} files
with syntax highlighting. It also offers productivity
shortcuts for frequently used LaTeX commands.

For example, the following commands are available by
default:

\begin{itemize}[noitemsep]
    \item \textbf{Ctrl+B}: Insert \verb|\textbf{}|.
    \item \textbf{Ctrl+I}: Insert \verb|\textit{}|.
    \item \textbf{Ctrl+U}: Insert \verb|\underline{}|.
    \item \textbf{Ctrl+Shift+E}: Insert \verb|\emph{}|.
    \item \textbf{Ctrl+Shift+M}: Insert inline mathematics.
\end{itemize}

When text is selected, formatting commands wrap the
selection. Otherwise, an empty command pair is inserted
and the caret is placed inside it.

All shortcuts can be customized through the standard
Eclipse key-binding preferences.

\subsection{Integrated PDF preview}

The PDF Preview view displays the compiled document
next to the source editor. It automatically follows the
active LaTeX file and loads the corresponding PDF.

The preview supports:
\begin{itemize}[noitemsep]
    \item Mouse-wheel scrolling.
    \item Zooming using Ctrl+mouse wheel or keyboard
    shortcuts.
    \item Fit-to-width and zoom commands.
    \item Highlighting of synchronized source locations.
\end{itemize}

PDF rendering is provided through Apache PDFBox
\cite{pdfbox}.

\subsection{Build integration}

EcTeX integrates the standard command-line tools of a TeX
distribution: a LaTeX engine (\texttt{pdflatex} by default, or
\texttt{xelatex} or \texttt{lualatex}) and a bibliography tool
(\texttt{bibtex} by default, or \texttt{biber}). Both are chosen
in the preferences.

The default build commands are:

\begin{center}
\begin{tabular}{ll}
\toprule
\textbf{Shortcut} & \textbf{Operation} \\
\midrule
Ctrl+F2 & Build PDF (one LaTeX pass) \\
Ctrl+F3 & Run the bibliography tool \\
Ctrl+F5 & Build All \\
\bottomrule
\end{tabular}
\end{center}

A LaTeX pass executes:

\begin{lstlisting}
pdflatex -interaction=nonstopmode \
         -file-line-error \
         -synctex=1
\end{lstlisting}

Compiler output is displayed in the EcTeX Build console.
LaTeX errors are also reported through Eclipse problem
markers.

For documents containing bibliographic references,
the recommended compilation sequence is:

\begin{enumerate}[noitemsep]
    \item Run LaTeX.
    \item Run the bibliography tool.
    \item Run LaTeX again.
    \item Run LaTeX once more.
\end{enumerate}

Build All performs exactly this sequence and stops at the
first error. For documents without a bibliography, the
preferences offer the shorter sequence of two LaTeX passes.

\section{Source and PDF Synchronization}

One of the central features of EcTeX is bidirectional
navigation between the LaTeX source and its compiled PDF.

This functionality is implemented using SyncTeX
\cite{synctex}.

\subsection{Forward synchronization}

Forward synchronization moves from the source editor
to the corresponding location in the PDF.

The user places the caret on a source line and presses
\texttt{Ctrl+Space}. The PDF preview scrolls to the
corresponding position and briefly highlights it.

\subsection{Inverse synchronization}

Inverse synchronization performs the opposite operation.
The user holds Ctrl and clicks on text in the PDF preview.
EcTeX then navigates to the corresponding source location.

This makes it easier to locate formatting problems,
review generated content and navigate longer documents.

\section{Technical Requirements}

EcTeX is implemented as an Eclipse plug-in and requires:

\begin{itemize}[noitemsep]
    \item Eclipse IDE with Java 17 or newer.
    \item A compatible TeX distribution, such as TeX Live
    or MiKTeX.
    \item The \texttt{pdflatex}, \texttt{bibtex} and
    \texttt{synctex} executables.
    \item Apache PDFBox 2.0.31 for PDF rendering.
\end{itemize}

If the required executables are not available through
the system \texttt{PATH}, their locations can be
configured through:

\begin{center}
\textit{Window $\rightarrow$ Preferences $\rightarrow$ EcTeX}
\end{center}

The plug-in has been built and install-tested against
Eclipse 2026-03.

\section{Known Limitations}

The current version focuses on providing a lightweight
and integrated LaTeX editing experience.

Some known limitations include:

\begin{itemize}[noitemsep]
    \item A project has one main file; to build another
    document of the same project, the main file has to be changed
    or cleared first.
    \item On high-DPI displays the PDF is rendered at the real
    resolution of the screen by an experimental setting that is
    on by default; it uses more memory and falls back to normal
    rendering if memory runs out.
    \item The PDF preview does not currently provide
    text-search functionality.
\end{itemize}

\bibliographystyle{plain}
\bibliography{references}

\end{document}
